\section{前向过程设计的数学验证}

\subsection{从一步到多步：闭式解的推导}

我们已经知道单步转移为：

$$x_t = \sqrt{\alpha_t} x_{t-1} + \sqrt{1 - \alpha_t} \epsilon_t$$

以 $t=2$ 为例，展开两步：

$$x_2 = \sqrt{\alpha_2} x_1 + \sqrt{1-\alpha_2} \epsilon_2$$

$$= \sqrt{\alpha_2}(\sqrt{\alpha_1} x_0 + \sqrt{1-\alpha_1} \epsilon_1) + \sqrt{1-\alpha_2} \epsilon_2$$

$$= \sqrt{\alpha_1 \alpha_2} x_0 + \sqrt{\alpha_2(1-\alpha_1)} \epsilon_1 + \sqrt{1-\alpha_2} \epsilon_2$$

由于 $\epsilon_1, \epsilon_2$ 独立同分布于 $\mathcal{N}(0, I)$：

$$\sqrt{\alpha_2(1-\alpha_1)} \epsilon_1 + \sqrt{1-\alpha_2} \epsilon_2 \sim \mathcal{N}(0, [\alpha_2(1-\alpha_1) + (1-\alpha_2)] I)$$

$$= \mathcal{N}(0, (1 - \alpha_1 \alpha_2) I)$$

因此 $q(x_2 | x_0) = \mathcal{N}(\sqrt{\alpha_1 \alpha_2} x_0, (1 - \alpha_1 \alpha_2) I) = \mathcal{N}(\sqrt{\bar{\alpha}_2} x_0, (1-\bar{\alpha}_2) I)$。

\begin{knowledgebox}{高斯分布的可加性}
如果 $a \sim \mathcal{N}(0, \sigma_1^2 I)$ 和 $b \sim \mathcal{N}(0, \sigma_2^2 I)$ 独立，则 $a + b \sim \mathcal{N}(0, (\sigma_1^2 + \sigma_2^2) I)$。这一性质是闭式解推导的核心工具。
\end{knowledgebox}

\subsection{验证极限行为}

当 $\bar{\alpha}_T \to 0$ 时：

$$q(x_T | x_0) = \mathcal{N}(\sqrt{\bar{\alpha}_T} x_0, (1-\bar{\alpha}_T)I) \to \mathcal{N}(0, I)$$

这验证了前向过程最终将数据完全转化为标准高斯噪声，与 ELBO 中先验匹配项的要求一致。

\subsection{本章小结}

通过逐步展开前向过程的递推关系并利用高斯分布的可加性，我们验证了闭式解的正确性，并确认了当 $T$ 足够大时，前向过程输出趋近标准高斯分布。


